% 第16批：扫描原书自测题5选择，共10题；题意独立转述。
\subsection{自测题5：气体动理论和热力学选择题}

\sourcetest{5.1.1}{115}{123}
\problemtext 封闭刚性容器中理想气体分子平均速率变为原来的2倍，问绝对温度和压强的倍率。
\workings 速率尺度 $\bar v\propto\sqrt T$，故 $T'=4T$；等容且粒子数不变，$p\propto T$，于是 $p'=4p$，选 \textbf{D}。
\crosscheck 题设速率是统计平均量；结论使用同种气体、理想气体、恒容前提。

\sourcetest{5.1.2}{115}{123}
\problemtext 活塞迅速外移而气体膨胀，判别气体是否经历平衡态以及能否以单条 $p$--$V$ 准静态曲线表述。
\workings 迅速过程通常存在非均匀压强和流动，系统内无单一瞬时热力学压强，非准静态平衡过程，不能以连续平衡态曲线准确表示，选 \textbf{C}。
\crosscheck 初态和充分弛豫后的末态可有明确 $p,V$，不等于中间全过程有平衡态曲线。

\sourcetest{5.1.3}{115}{123}
\problemtext 相同容积容器中的氧气与氦气初温初压相同；给氧气5J使其升高某温差，氦气升高同温差需多少热量？
\workings 同温同压同体积意味着粒子数相同。常温氧气刚性双原子，$C_V=5R/2$；氦气单原子，$C_V=3R/2$。等容下 $\boxed{Q_{\rm He}/Q_{\rm O_2}=3/5}$，故 $Q_{\rm He}=3$J，选 \textbf{C}。
\crosscheck 按理想气体、忽略振动态激发的常温近似；容器均刚性。

\sourcetest{5.1.4}{115}{123}
\problemtext 题5.1.2图中两个卡诺循环共用高低等温线，但第二个循环面积较大；比较净功与效率。
\workings 一个循环的净功是 $pV$ 图中围成的面积；第二循环面积变大故净功增加。卡诺效率仅与热源温度相关，$\eta=1-T_2/T_1$ 不变，选 \textbf{D}。
\crosscheck 题图改变循环体积尺度，未改变冷热源温度。

\sourcetest{5.1.5}{115}{123}
\problemtext 在温度为327°C和27°C的热源之间工作的理想可逆热机，最大效率是多少？
\workings $T_h=600.15$K、$T_c=300.15$K，$\eta_{\max}=1-T_c/T_h\approx\boxed{50\%}$，选 \textbf{B}。
\crosscheck 应换算成开尔文后才取比；按通常近似273K换算也得到50%。

\sourcetest{5.1.6}{115--116}{NaN}
\problemtext 理想气体从 $(V_0,T_0)$ 先绝热膨胀到 $2V_0$，等容升温回 $T_0$，再等温压缩回初态；判断循环净功。
\workings 令 $\gamma>1$。绝热末温 $T_a=T_0\,2^{1-\gamma}$，绝热功 $W_a=nRT_0[1-2^{1-\gamma}]/(\gamma-1)$；等容功0，等温压缩功 $W_c=-nRT_0\ln2$。因 $1-e^{-z}<z$（$z=(\gamma-1)\ln2$），$W_a+W_c<0$，选 \textbf{B}。
\crosscheck 一循环 $\Delta U=0$，净功为负意味着需净输入功。

\sourcetest{5.1.7}{116}{124}
\problemtext 一定量理想气体在真空中绝热自由膨胀，体积增大；问内能、熵变化。
\workings 无热量交换 $Q=0$，且向真空不作体积功 $W=0$，故 $\Delta U=0$、$T$ 不变；熵变 $\Delta S=nR\ln(V_2/V_1)>0$，选 \textbf{A}。
\crosscheck 过程不可逆，不能把实际膨胀误判为可逆绝热等熵过程。

\sourcetest{5.1.8}{116}{124}
\problemtext 理想气体只接触一个热源作等温膨胀，吸热全部转化为对外做功；这是否违反热力学定律？
\workings 对理想气体等温 $\Delta U=0$，由第一定律 $Q=W>0$；热力学第二定律的开尔文表述禁止的是\emph{循环}唯一效果，不禁止单次等温膨胀。两定律均不违反，选 \textbf{C}。
\crosscheck 若要回到初态必须另行进行复原步骤，不能将此过程声称为单热源循环热机。

\sourcetest{5.1.9}{116}{124}
\problemtext 刚性双原子氮气经可逆绝热压缩，压强增至2倍，求分子平均速率的倍率。
\workings $\gamma=7/5$，$T_2/T_1=(p_2/p_1)^{(\gamma-1)/\gamma}=2^{2/7}$；$\bar v_2/\bar v_1=\sqrt{T_2/T_1}=\boxed{2^{1/7}}$，选 \textbf{D}。
\crosscheck 若问分子均方根速率、平均速率或最概然速率，固定分子质量下倍率均为 $\sqrt T$。

\sourcetest{5.1.10}{116}{124}
\problemtext 理想气体绝热向真空自由膨胀，体积翻倍；比较初末温度与平均自由程。
\workings 自由膨胀理想气体 $\Delta U=0$，故 $T_1=T_2$；气体数密度减半，固定温度时 $\lambda\propto1/n$，故 $\lambda_2=2\lambda_1$。选 \textbf{B}。
\crosscheck 需要同种气体的有效碰撞直径不变；与可逆绝热降温不同。
